\pagestyle{empty}
\tight
\begin{tblr}{width=\textwidth,colspec={|Q[c,m,1.9cm]|X[l,m]|},hlines,vlines,hspan=minimal,row{1}={bg=headblue},rowsep=1pt,colsep=3pt}
\SetCell{c,m}\hfil\logo\hfil & \textbf{POLITEKNIK NEGERI MALANG}\par\textbf{JURUSAN TEKNOLOGI INFORMASI}\par\textbf{PROGRAM STUDI: D4 TEKNIK INFORMATIKA} \\
\SetCell[c=2]{c} \textbf{RENCANA PEMBELAJARAN SEMESTER (RPS)} & \\
\end{tblr}
\begin{longtblr}{width=\textwidth,colspec={|Q[l,m,1.9cm]|X[0.85,l,m]|X[1.45,l,m]|X[0.75,l,m]|X[1.15,l,m]|},hlines,vlines,hspan=minimal,rowsep=1pt,colsep=2pt,row{1,3}={bg=headgray,font=\bfseries}}
MATA KULIAH & KODE & BOBOT (SKS)/jam & SEMESTER & TGL. PENYUSUNAN \\
Teknologi Terapan & RTI256207 & 2 SKS/4 jam & 6 & 1 Juli 2025 \\
OTORISASI & Dosen Pengembang RPS & Koordinator RMK & \SetCell[c=2]{l} Ka PRODI &  \\
 & Triana Fatmawati, S.T., M.T. & Triana Fatmawati, S.T., M.T. & \SetCell[c=2]{l}{\vspace{8mm}\par Dr. Ely Setyo Astuti, S.T., M.T.} &  \\
\SetCell[r=12]{l} Capaian Pembelajaran (CP) & \SetCell[c=4]{l,bg=headgray,font=\bfseries} Capaian Pembelajaran Lulusan Program Studi (CPL-Prodi) &  &  &  \\
 & CPL02 & \SetCell[c=3]{l} Mampu menerapkan prinsip rekayasa sistem dan perangkat lunak untuk menghasilkan solusi aplikasi multi-platform sesuai kebutuhan. &  &  \\
 & CPL05 & \SetCell[c=3]{l} Mampu menghasilkan solusi teknologi komputasi dan infrastruktur digital yang sesuai dengan kebutuhan organisasi dan pengguna. &  &  \\
 & CPL08 & \SetCell[c=3]{l} Mampu mengelola proyek teknologi informasi secara profesional melalui perencanaan, koordinasi, komunikasi, dan optimalisasi sumber daya. &  &  \\
 & \SetCell[c=4]{l,bg=headgray,font=\bfseries} Capaian Pembelajaran mata kuliah (CPL-MK) &  &  &  \\
 & CPMK0204 & \SetCell[c=3]{l} Mahasiswa mampu menganalisis kebutuhan dan merancang arsitektur solusi teknologi terapan berbasis kebutuhan nyata. &  &  \\
 & CPMK0503 & \SetCell[c=3]{l} Mahasiswa mampu mengimplementasikan solusi teknologi terapan yang efisien dan memenuhi kebutuhan operasional. &  &  \\
 & CPMK0801 & \SetCell[c=3]{l} Mahasiswa mampu merencanakan dan mengelola implementasi teknologi terapan di lingkungan industri. &  &  \\
 & \SetCell[c=4]{l,bg=headgray,font=\bfseries} Kemampuan akhir tiap tahapan belajar (Sub-CPMK) &  &  &  \\
 & SCPMK0204-05601 & \SetCell[c=3]{l} Mahasiswa mampu melakukan analisis kebutuhan teknologi industri dan merancang solusi teknologi terapan yang tepat guna. &  &  \\
 & SCPMK0503-05602 & \SetCell[c=3]{l} Mahasiswa mampu mengimplementasikan solusi teknologi terapan mulai dari prototipe hingga deployment produksi. &  &  \\
 & SCPMK0801-05603 & \SetCell[c=3]{l} Mahasiswa mampu menyusun rencana implementasi teknologi, mengelola risiko, dan melaporkan kemajuan kepada stakeholder. &  &  \\
\SetCell[r=2]{l} Penilaian & \SetCell[c=4]{l} *Nilai CPMK didapat dari agregasi nilai SUB-CPMK yang dibebankan &  &  &  \\
 & \SetCell[c=4]{l}{\tiny\begin{tblr}{width=\linewidth,colspec={|Q[c,m,0.1\linewidth]|Q[c,m,0.16\linewidth]|X[l,m]|X[l,m]|X[l,m]|X[l,m]|X[l,m]|X[l,m]|Q[c,m,0.08\linewidth]|},hlines,vlines,hspan=minimal,rowsep=1pt,colsep=0.7pt,row{1,2}={bg=headgray,font=\bfseries}}
Kode CPMK & Kode SCPMK & \SetCell[c=6]{c} Bobot per Bentuk Penilaian & & & & & & Total Bobot \\
 & & Analisis & Sprint 1-2 & UTS & Sprint 3-4 & Laporan & UAS Presentasi & \\
CPMK0204 & SCPMK0204-05601 & 8\%: analisis kebutuhan teknologi & 8\%: sprint 1-2 analisis & 0\% & 8\%: sprint 3-4 analisis & 5\%: laporan analisis & 0\% & 29\% \\
CPMK0503 & SCPMK0503-05602 & 0\% & 12\%: implementasi sprint 1-2 & 0\% & 12\%: implementasi sprint 3-4 & 5\%: laporan implementasi & 10\%: demo produk & 39\% \\
CPMK0801 & SCPMK0801-05603 & 7\%: rencana implementasi & 0\% & 15\%: milestone review & 0\% & 0\% & 10\%: presentasi stakeholder & 32\% \\
\SetCell[c=8]{r}\textbf{Total Per Penilaian} & & & & & & & & \textbf{100\%} \\
\end{tblr}} &  &  &  \\
Diskripsi Singkat Mata Kuliah & \SetCell[c=4]{l} Teknologi Terapan adalah capstone jalur magang/industri yang mengintegrasikan analisis kebutuhan industri, perancangan, dan implementasi solusi teknologi tepat guna dalam konteks industri nyata melalui pendekatan iteratif dan kolaboratif bersama stakeholder. &  &  &  \\
Bahan Kajian / Materi Pembelajaran & \SetCell[c=4]{l}\begin{enumerate}\item Needs assessment, technology mapping, dan feasibility study untuk solusi teknologi industri\item Perencanaan implementasi teknologi: roadmap, manajemen risiko, dan stakeholder management\item Implementasi solusi iteratif: prototipe, pengembangan, integrasi, dan optimasi\item Deployment produksi, transfer knowledge, dan evaluasi dampak teknologi\end{enumerate} &  &  &  \\
\SetCell[r=2]{l} Pustaka & \SetCell[c=4]{l} \textbf{Utama:}\begin{enumerate}\item Rogers, E. 2003. Diffusion of Innovations. Free Press.\item Christensen, C. 2016. The Innovator's Dilemma. HBR Press.\item Brown, T. 2009. Change by Design. HarperBusiness.\end{enumerate} &  &  &  \\
 & \SetCell[c=4]{l} \textbf{Pendukung:} Materi LMS, dokumentasi resmi perangkat lunak, dan jobsheet praktikum. &  &  &  \\
Media Pembelajaran & \SetCell[c=2]{l} \textbf{Software:} IDE, prototyping tools, cloud platform, project management tools, LMS. &  & \SetCell[c=2]{l} \textbf{Hardware:} Komputer/laptop, proyektor. &  \\
Nama Dosen Pengampu & \SetCell[c=4]{l} Tim Pengajar Program Studi D4 TI &  &  &  \\
Matakuliah Syarat & \SetCell[c=4]{l} - &  &  &  \\
\end{longtblr}
\begin{longtblr}{width=\textwidth,colspec={|Q[c,m,0.06\textwidth]|X[1.35,l,m]|X[1.08,l,m]|X[1.16,l,m]|X[0.7,l,m]|X[1.24,l,m]|X[1.06,l,m]|X[0.98,l,m]|Q[c,m,0.05\textwidth]|},hlines,vlines,hspan=minimal,rowhead=2,row{1,2}={bg=headgreen,font=\bfseries},rowsep=1pt,colsep=1.5pt}
Mgg.\par Ke & Kemampuan Akhir Yang Direncanakan (Sub-CP-MK) & Bahan kajian (Materi Pembelajaran) & Bentuk dan Metode Pembelajaran & Estimasi Waktu & Pengalaman Belajar Mahasiswa & Kriteria \& Bentuk Penilaian & Indikator Penilaian & Bobot Penilaian (\%) \\
(1) & (2) & (3) & (4) & (5) & (6) & (7) & (8) & (9) \\
1 & SCPMK0204-05601 & Orientasi teknologi terapan dan identifikasi domain. & Modalitas: Blended Learning. Bentuk: Luring/Daring. Metode: Project Based Learning, iterative prototyping, stakeholder engagement, dan refleksi. & 1 x 6 x 50' workshop industri; 1 x 6 x 50' dokumentasi dan tugas mandiri. & Mahasiswa mengerjakan aktivitas untuk orientasi teknologi terapan dan identifikasi domain serta menunjukkan evidence hasil belajar. & Bentuk penilaian: Orientasi Teknologi Terapan. Mengacu rubrik RTM. & Ketepatan analisis; kualitas solusi; validasi dan dokumentasi. & 3 \\
2 & SCPMK0204-05601 & Needs assessment dan technology mapping. & Modalitas: Blended Learning. Bentuk: Luring/Daring. Metode: Project Based Learning, iterative prototyping, stakeholder engagement, dan refleksi. & 1 x 6 x 50' workshop industri; 1 x 6 x 50' dokumentasi dan tugas mandiri. & Mahasiswa mengerjakan aktivitas untuk needs assessment dan technology mapping serta menunjukkan evidence hasil belajar. & Bentuk penilaian: Needs Assessment dan Technology Mapping. Mengacu rubrik RTM. & Ketepatan analisis; kualitas solusi; validasi dan dokumentasi. & 3 \\
3 & SCPMK0204-05601 & Feasibility study dan pemilihan teknologi. & Modalitas: Blended Learning. Bentuk: Luring/Daring. Metode: Project Based Learning, iterative prototyping, stakeholder engagement, dan refleksi. & 1 x 6 x 50' workshop industri; 1 x 6 x 50' dokumentasi dan tugas mandiri. & Mahasiswa mengerjakan aktivitas untuk feasibility study dan pemilihan teknologi serta menunjukkan evidence hasil belajar. & Bentuk penilaian: Feasibility Study dan Pemilihan Teknologi. Mengacu rubrik RTM. & Ketepatan analisis; kualitas solusi; validasi dan dokumentasi. & 4 \\
4 & SCPMK0801-05603 & Perencanaan implementasi: roadmap dan manajemen risiko. & Modalitas: Blended Learning. Bentuk: Luring/Daring. Metode: Project Based Learning, iterative prototyping, stakeholder engagement, dan refleksi. & 1 x 6 x 50' workshop industri; 1 x 6 x 50' dokumentasi dan tugas mandiri. & Mahasiswa mengerjakan aktivitas untuk perencanaan implementasi roadmap dan manajemen risiko serta menunjukkan evidence hasil belajar. & Bentuk penilaian: Perencanaan Implementasi. Mengacu rubrik RTM. & Ketepatan analisis; kualitas solusi; validasi dan dokumentasi. & 4 \\
5 & SCPMK0503-05602 & Sprint 1: setup infrastruktur dan prototipe awal. & Modalitas: Blended Learning. Bentuk: Luring/Daring. Metode: Project Based Learning, iterative prototyping, stakeholder engagement, dan refleksi. & 1 x 6 x 50' workshop industri; 1 x 6 x 50' dokumentasi dan tugas mandiri. & Mahasiswa mengerjakan aktivitas sprint 1 untuk setup infrastruktur dan prototipe awal serta menunjukkan evidence hasil belajar. & Bentuk penilaian: Sprint 1 Setup dan Prototipe Awal. Mengacu rubrik RTM. & Ketepatan analisis; kualitas solusi; validasi dan dokumentasi. & 5 \\
6 & SCPMK0503-05602 & Sprint 1: validasi prototipe dengan pengguna. & Modalitas: Blended Learning. Bentuk: Luring/Daring. Metode: Project Based Learning, iterative prototyping, stakeholder engagement, dan refleksi. & 1 x 6 x 50' workshop industri; 1 x 6 x 50' dokumentasi dan tugas mandiri. & Mahasiswa mengerjakan aktivitas sprint 1 untuk validasi prototipe dengan pengguna serta menunjukkan evidence hasil belajar. & Bentuk penilaian: Sprint 1 Validasi Prototipe. Mengacu rubrik RTM. & Ketepatan analisis; kualitas solusi; validasi dan dokumentasi. & 5 \\
7 & SCPMK0801-05603 & Sprint 1 review dan stakeholder feedback. & Modalitas: Blended Learning. Bentuk: Luring/Daring. Metode: Project Based Learning, iterative prototyping, stakeholder engagement, dan refleksi. & 1 x 6 x 50' workshop industri; 1 x 6 x 50' dokumentasi dan tugas mandiri. & Mahasiswa mengerjakan aktivitas sprint 1 review dan stakeholder feedback serta menunjukkan evidence hasil belajar. & Bentuk penilaian: Sprint 1 Review dan Stakeholder Feedback. Mengacu rubrik RTM. & Ketepatan analisis; kualitas solusi; validasi dan dokumentasi. & 5 \\
8 & SCPMK0801-05603 & UTS Milestone: evaluasi prototipe dan rencana lanjutan. & Modalitas: Blended Learning. Bentuk: Luring/Daring. Metode: Project Based Learning, iterative prototyping, stakeholder engagement, dan refleksi. & 1 x 6 x 50' workshop industri; 1 x 6 x 50' dokumentasi dan tugas mandiri. & Mahasiswa mengerjakan aktivitas UTS berupa evaluasi prototipe dan rencana lanjutan serta menunjukkan evidence hasil belajar. & Bentuk penilaian: UTS Milestone Evaluasi Prototipe. Mengacu rubrik RTM. & Ketepatan analisis; kualitas solusi; validasi dan dokumentasi. & 15 \\
9 & SCPMK0503-05602 & Sprint 2: pengembangan solusi lanjutan. & Modalitas: Blended Learning. Bentuk: Luring/Daring. Metode: Project Based Learning, iterative prototyping, stakeholder engagement, dan refleksi. & 1 x 6 x 50' workshop industri; 1 x 6 x 50' dokumentasi dan tugas mandiri. & Mahasiswa mengerjakan aktivitas sprint 2 untuk pengembangan solusi lanjutan serta menunjukkan evidence hasil belajar. & Bentuk penilaian: Sprint 2 Pengembangan Solusi Lanjutan. Mengacu rubrik RTM. & Ketepatan analisis; kualitas solusi; validasi dan dokumentasi. & 4 \\
10 & SCPMK0503-05602 & Sprint 2: integrasi dengan sistem yang ada. & Modalitas: Blended Learning. Bentuk: Luring/Daring. Metode: Project Based Learning, iterative prototyping, stakeholder engagement, dan refleksi. & 1 x 6 x 50' workshop industri; 1 x 6 x 50' dokumentasi dan tugas mandiri. & Mahasiswa mengerjakan aktivitas sprint 2 untuk integrasi dengan sistem yang ada serta menunjukkan evidence hasil belajar. & Bentuk penilaian: Sprint 2 Integrasi Sistem. Mengacu rubrik RTM. & Ketepatan analisis; kualitas solusi; validasi dan dokumentasi. & 5 \\
11 & SCPMK0801-05603 & Sprint 2 review dan Sprint 3 planning. & Modalitas: Blended Learning. Bentuk: Luring/Daring. Metode: Project Based Learning, iterative prototyping, stakeholder engagement, dan refleksi. & 1 x 6 x 50' workshop industri; 1 x 6 x 50' dokumentasi dan tugas mandiri. & Mahasiswa mengerjakan aktivitas sprint 2 review dan sprint 3 planning serta menunjukkan evidence hasil belajar. & Bentuk penilaian: Sprint 2 Review Sprint 3 Planning. Mengacu rubrik RTM. & Ketepatan analisis; kualitas solusi; validasi dan dokumentasi. & 4 \\
12 & SCPMK0503-05602 & Sprint 3: optimasi dan stress testing. & Modalitas: Blended Learning. Bentuk: Luring/Daring. Metode: Project Based Learning, iterative prototyping, stakeholder engagement, dan refleksi. & 1 x 6 x 50' workshop industri; 1 x 6 x 50' dokumentasi dan tugas mandiri. & Mahasiswa mengerjakan aktivitas sprint 3 untuk optimasi dan stress testing serta menunjukkan evidence hasil belajar. & Bentuk penilaian: Sprint 3 Optimasi dan Stress Testing. Mengacu rubrik RTM. & Ketepatan analisis; kualitas solusi; validasi dan dokumentasi. & 5 \\
13 & SCPMK0204-05601 & Sprint 3: penyesuaian berbasis feedback industri. & Modalitas: Blended Learning. Bentuk: Luring/Daring. Metode: Project Based Learning, iterative prototyping, stakeholder engagement, dan refleksi. & 1 x 6 x 50' workshop industri; 1 x 6 x 50' dokumentasi dan tugas mandiri. & Mahasiswa mengerjakan aktivitas sprint 3 untuk penyesuaian berbasis feedback industri serta menunjukkan evidence hasil belajar. & Bentuk penilaian: Sprint 3 Penyesuaian Feedback Industri. Mengacu rubrik RTM. & Ketepatan analisis; kualitas solusi; validasi dan dokumentasi. & 4 \\
14 & SCPMK0503-05602 & Sprint 4: deployment produksi dan transfer knowledge. & Modalitas: Blended Learning. Bentuk: Luring/Daring. Metode: Project Based Learning, iterative prototyping, stakeholder engagement, dan refleksi. & 1 x 6 x 50' workshop industri; 1 x 6 x 50' dokumentasi dan tugas mandiri. & Mahasiswa mengerjakan aktivitas sprint 4 untuk deployment produksi dan transfer knowledge serta menunjukkan evidence hasil belajar. & Bentuk penilaian: Sprint 4 Deployment dan Transfer Knowledge. Mengacu rubrik RTM. & Ketepatan analisis; kualitas solusi; validasi dan dokumentasi. & 5 \\
15 & SCPMK0503-05602 & Dokumentasi teknis dan laporan implementasi. & Modalitas: Blended Learning. Bentuk: Luring/Daring. Metode: Project Based Learning, iterative prototyping, stakeholder engagement, dan refleksi. & 1 x 6 x 50' workshop industri; 1 x 6 x 50' dokumentasi dan tugas mandiri. & Mahasiswa mengerjakan aktivitas untuk dokumentasi teknis dan laporan implementasi serta menunjukkan evidence hasil belajar. & Bentuk penilaian: Dokumentasi Teknis dan Laporan Implementasi. Mengacu rubrik RTM. & Ketepatan analisis; kualitas solusi; validasi dan dokumentasi. & 5 \\
16 & SCPMK0801-05603 & Persiapan presentasi final kepada stakeholder. & Modalitas: Blended Learning. Bentuk: Luring/Daring. Metode: Project Based Learning, iterative prototyping, stakeholder engagement, dan refleksi. & 1 x 6 x 50' workshop industri; 1 x 6 x 50' dokumentasi dan tugas mandiri. & Mahasiswa mengerjakan aktivitas untuk persiapan presentasi final kepada stakeholder serta menunjukkan evidence hasil belajar. & Bentuk penilaian: Persiapan Presentasi Final. Mengacu rubrik RTM. & Ketepatan analisis; kualitas solusi; validasi dan dokumentasi. & 4 \\
17 & SCPMK0801-05603, SCPMK0503-05602 & UAS: presentasi teknologi terapan dan demo kepada industri. & Modalitas: Blended Learning. Bentuk: Luring/Daring. Metode: Project Based Learning, iterative prototyping, stakeholder engagement, dan refleksi. & 1 x 6 x 50' workshop industri; 1 x 6 x 50' dokumentasi dan tugas mandiri. & Mahasiswa mengerjakan aktivitas UAS berupa presentasi teknologi terapan dan demo kepada industri serta menunjukkan evidence hasil belajar. & Bentuk penilaian: UAS Presentasi Teknologi Terapan dan Demo. Mengacu rubrik RTM. & Ketepatan analisis; kualitas solusi; validasi dan dokumentasi. & 20 \\
\end{longtblr}
